\documentclass[a4paper,10.5pt]{article}
\usepackage[utf8]{inputenc}
\usepackage[T1]{fontenc}
\usepackage[french]{babel}
\usepackage[left=1cm,right=1cm,top=.5cm,bottom=0cm]{geometry}
\usepackage{enumitem}
\usepackage{hyperref}
\usepackage{xcolor}
\usepackage{titlesec}
\usepackage{fontawesome5}
\usepackage{lmodern}
\usepackage[normalem]{ulem}

% --- Couleurs Thales ---
\definecolor{primary}{RGB}{0, 56, 101}
\definecolor{secondary}{RGB}{80, 80, 80}

% --- Configuration des liens ---
\hypersetup{
    colorlinks=true,
    linkcolor=primary,
    filecolor=primary,
    urlcolor=primary,
}

% --- Formatage des titres ---
\titleformat{\section}{\large\bfseries\color{primary}\uppercase}{}{0em}{}[\titlerule]
\titlespacing{\section}{0pt}{8pt}{5pt}

% --- Commandes personnalisées ---
\newcommand{\resumeItem}[1]{\item\small{#1}}
\newcommand{\resumeSubheading}[4]{
  \vspace{-1pt}\item
    \begin{tabular*}{0.97\textwidth}[t]{l@{\extracolsep{\fill}}r}
      \textbf{#1} & \textit{\small #2} \\
      \textit{\small #3} & \textit{\small #4} \\
    \end{tabular*}\vspace{-5pt}
}

\begin{document}

% --- EN-TÊTE ---
\begin{center}
    {\Large \textbf{Loïc Aron MBASSI EWOLO}} \\ \vspace{4pt}
    \textbf{\large Stage Ingénieur Développement FPGA (MPSoC/RFSoC)} \\
    \textit{\small Disponible dès \textbf{avril 2026} pour 4 à 6 mois} \\ \vspace{4pt}
    \small
    \faMapMarker\ Cergy, France (Mobile France entière) \quad
    \faPhone\ (+33) 7 59 52 33 98 \quad
    \faEnvelope\ \href{mailto:loicaron.mbassiewolo@ensea.fr}{loicaron.mbassiewolo@ensea.fr} \\ \vspace{2pt}
    \faLinkedin\ \href{https://www.linkedin.com/in/loic-aron-mbassi-ewolo-3942a9246/}{linkedin.com/in/loic-mbassi} \quad
    \faGithub\ \href{https://github.com/Nameless0l}{github.com/Nameless0l} \quad
    \faGlobe\ \href{https://mbassiloic.tech}{mbassiloic.tech}
\end{center}

\vspace{-6pt}

% --- PROFIL ---
\noindent\small{Élève-ingénieur à l'\textbf{ENSEA} en double diplôme, spécialisé en \textbf{Systèmes Embarqués} et Traitement du Signal. Expérience en développement \textbf{VHDL}, programmation \textbf{C bas niveau} et conception hardware. Passionné par l'électronique numérique et le développement sur cibles FPGA. Je recherche un stage stimulant chez \textbf{Thales} pour contribuer au développement de solutions FPGA de nouvelle génération.}

% --- FORMATION ---
\section{Formation}
\begin{itemize}[leftmargin=*, label=\tiny$\bullet$, nosep]
    \item \textbf{ENSEA - École Nationale Supérieure de l'Électronique et de ses Applications} \hfill \textit{Cergy | 2025 -- Présent} \\
    \small{Ingénieur 2A, Double Diplôme -- \textbf{Systèmes Embarqués}, Traitement du Signal, option Drones.} \\
    \textit{\small{Cours clés : Conception Numérique FPGA (\textbf{VHDL}), Architecture des Processeurs, Traitement du Signal, Automatique.}}
    \vspace{2pt}
    \item \textbf{École Polytechnique de Yaoundé} (1\textsuperscript{ère} école d'ingénieurs CEMAC) \hfill \textit{Cameroun | 2021 -- 2025} \\
    \small{Diplôme d'Ingénieur de Conception (Master 2) : \textbf{Mathématiques Appliquées}, Génie Logiciel, Algorithmique.} \\
    \small{Classement : \textbf{20\textsuperscript{ème}/108} -- Moyenne : 14,03/20.}
    \vspace{2pt}
    \item \textbf{Université de Yaoundé I} \hfill \textit{Cameroun | 2020 -- 2022} \\
    \small{Licence Informatique \& Mathématiques : Algorithmique \textbf{C}, Analyse Numérique, Algèbre Linéaire.}
\end{itemize}

% --- COMPÉTENCES TECHNIQUES ---
\section{Compétences Techniques}
\begin{itemize}[leftmargin=*, nosep]
    \item \small{\textbf{FPGA \& HDL :} \textbf{VHDL}, Simulation (ModelSim/QuestaSim), Synthèse (Vivado), Conception numérique.}
    \item \small{\textbf{Embarqué \& Hardware :} \textbf{C/C++} (bas niveau), STM32, Arduino, Raspberry Pi, Linux Embarqué, I2C, SPI, UART.}
    \item \small{\textbf{Outils \& Environnements :} Git, Docker, \textbf{MATLAB/Simulink}, Python, Altium/KiCad (notions PCB).}
    \item \small{\textbf{Méthodologie :} Tests unitaires, Intégration Continue, Documentation technique, Travail en équipe Agile.}
\end{itemize}

% --- PROJETS ACADÉMIQUES FPGA/EMBARQUÉ ---
\section{Projets Académiques \& Techniques}
\begin{itemize}[leftmargin=*, label={-}]

\item \textbf{Conception Numérique sur FPGA (VHDL)} \hfill \textit{Projet ENSEA | 2025 -- En cours}
\vspace{-2pt}
    \begin{itemize}[leftmargin=0.4cm, nosep, label=\tiny$\bullet$]
        \item \small{Développement de descriptions matérielles en \textbf{VHDL} pour circuits logiques numériques.}
        \item \small{Simulation comportementale et validation fonctionnelle sur carte \textbf{FPGA}.}
        \item \small{Synthèse et implémentation via \textbf{Vivado}, analyse des ressources et timing.}
    \end{itemize}
\vspace{3pt}

\item \textbf{Fault Detection \& Fault-Tolerant Control -- Drone} \hfill \textit{Projet d'Option ENSEA | 2026}
\vspace{-2pt}
    \begin{itemize}[leftmargin=0.4cm, nosep, label=\tiny$\bullet$]
        \item \small{Modélisation dynamique d'un drone en conditions nominales et dégradées (perte actionneurs, biais capteurs).}
        \item \small{Conception d'un schéma de \textbf{détection de défauts} basé sur observateurs (résidus, équations de parité).}
        \item \small{Implémentation et simulations sous \textbf{MATLAB/Simulink}.}
    \end{itemize}
\vspace{3pt}

\item \textbf{Challenge CoHoMa -- Robotique Autonome Militaire} \hfill \textit{ENSEA/Armée Française | 2025-2026}
\vspace{-2pt}
    \begin{itemize}[leftmargin=0.4cm, nosep, label=\tiny$\bullet$]
        \item \small{Optimisation d'algorithmes de \textbf{SLAM} et déploiement IA embarquée sur \textbf{Nvidia Jetson}.}
        \item \small{Intégration capteurs (Lidar VLP16, caméras profondeur), gestion flux temps réel.}
        \item \small{Environnement \textbf{Docker} pour isolation et reproductibilité, programmation \textbf{C/C++}.}
    \end{itemize}
\vspace{3pt}

\item \textbf{PROJET-TAXI -- Simulation Système Temps Réel} \hfill \textit{Projet Académique Polytechnique | 2024}
\vspace{-2pt}
    \begin{itemize}[leftmargin=0.4cm, nosep, label=\tiny$\bullet$]
        \item \small{Développement en \textbf{C} d'un noyau de simulation multiprocessus pour flotte de véhicules.}
        \item \small{Gestion \textbf{mémoire partagée} (SHM), communication inter-processus (IPC), synchronisation par \textbf{signaux UNIX}.}
        \item \small{Implémentation d'un ordonnanceur \textbf{FIFO} avec prévention des interblocages.}
    \end{itemize}
\vspace{3pt}

\item \textbf{Harmony Gloves -- Gants Instrumentés pour Langue des Signes} \hfill \textit{Laboratoire IOTA | 2024}
\vspace{-2pt}
    \begin{itemize}[leftmargin=0.4cm, nosep, label=\tiny$\bullet$]
        \item \small{Prototypage hardware : conception circuit, \textbf{soudure} composants (capteurs flex, IMU).}
        \item \small{Programmation microcontrôleur (\textbf{STM32/Arduino}), protocoles \textbf{I2C/SPI}.}
        \item \small{Fusion de données capteurs pour classification temps réel (TinyML).}
    \end{itemize}

\end{itemize}

% --- EXPÉRIENCE PROFESSIONNELLE ---
\section{Expérience Professionnelle}
\begin{itemize}[leftmargin=*, label={-}]

    \resumeSubheading
      {Stage Recherche -- Laboratoire IoT (TinyML)}{Oct. 2023 -- Jan. 2024}
      {École Polytechnique de Yaoundé}{Yaoundé, Cameroun}
      \begin{itemize}[leftmargin=0.4cm, nosep, label=\tiny$\bullet$]
        \item \small{Optimisation de modèles ML pour \textbf{hardware limité} (contraintes mémoire, temps réel).}
        \item \small{Déploiement sur \textbf{Arduino/Raspberry Pi} avec TensorFlow Lite, programmation \textbf{C}.}
      \end{itemize}
    \vspace{2pt}

    \resumeSubheading
      {Assistant de Recherche -- Interprétabilité IA}{Juin 2025 -- Sept. 2025}
      {MILA (Institut Québécois d'IA)}{Québec, Canada (Remote)}
      \begin{itemize}[leftmargin=0.4cm, nosep, label=\tiny$\bullet$]
        \item \small{Recherche sur la généralisation tardive (Grokking) des réseaux de neurones.}
        \item \small{Implémentation de pipelines de tests sous \textbf{PyTorch}, analyse mathématique de convergence.}
      \end{itemize}

\end{itemize}

% --- DISTINCTIONS & CERTIFICATIONS ---
\section{Distinctions \& Soft Skills}
\begin{itemize}[leftmargin=*, nosep]
    \item \textbf{Hackathons :} 1\textsuperscript{er} Prix IndabaX Africa 2025 | 1\textsuperscript{er} CITS National | 1\textsuperscript{er} Mountain Hub 2024.
    \item \textbf{Leadership :} Président Club Informatique, Chef Cellule Projet, Lead Cellule IA (Polytechnique Yaoundé).
    \item \textbf{Langues :} Français (natif), Anglais (professionnel/technique).
    \item \textbf{Qualités :} Rigueur, autonomie, esprit d'analyse, travail en équipe, curiosité technique.
\end{itemize}

\end{document}
